\begin{table}[htbp]
\centering
\caption{Border-vs-Interior Estimates: Cannabis Revenue per Dispensary Store (CO + WA)}
\label{tab:revenue_perstore_border70}
\begin{tabular}{lcccc}
\toprule
Specification & Coefficient & SE (HC1) & SE (county-cl.) & $p$-value (county-cl.) \\
\midrule
No controls & $+0.502$ & $0.101$ & $0.233$ & $0.031$ \\
+ log(pop) + year FE + state FE & $+0.395$ & $0.074$ & $0.163$ & $0.015$ \\
\quad excluding Sedgwick County, CO & $+0.302$ & $0.068$ & $0.147$ & $0.040$ \\
\midrule
\multicolumn{5}{l}{\textit{Single-state estimates, + log(pop) + year FE (both HC1 $p<0.001$)}} \\
Colorado only (42 counties) & $+0.552$ & $<0.001$ & $0.199$ & $0.005$ \\
Washington only (29 counties) & $+0.441$ & $<0.001$ & $0.263$ & $0.077$ \\
\bottomrule
\end{tabular}
\begin{minipage}{0.95\textwidth}
\vspace{0.3em}
\footnotesize
\textit{Notes:} Outcome is log(total county cannabis revenue $\div$ number of currently-open dispensary
stores). Store counts are cumulative-open-by-year, built from geocoded dispensary license records
point-in-polygon assigned to counties, counted as open if the license's open date precedes each
fiscal/calendar year's own reporting cutoff (matching each state's revenue-reporting convention). Sample
restricted to county-years with at least one open store, $n=706$ (71 counties, 10 border and 61
interior). Unlike the revenue-per-capita and dispensary-density checks reported elsewhere in this
project -- both of which are null -- revenue per store is significantly \emph{higher} in border counties,
robust to dropping Sedgwick County (the single-county outlier that drove an earlier, non-robust
50-mile-threshold result) and to estimating each state separately. This is the only significant result
of the checks reported in this project's border-competition research thread, and it is directly
consistent with the license-cap explanation proposed when the original entry-density check came back null:
if border-county demand is genuinely elevated but the number of licensed stores is capped by state
regulation, the excess demand should appear as higher revenue per existing store rather than as more
stores or higher aggregate per-capita revenue -- exactly the pattern observed here.
\end{minipage}
\end{table}
